\documentclass[11pt]{article}
% Shared preamble for the qurl-theory reports.
% Build either report from this folder with:  latexmk -pdf <file>.tex
% (or: pdflatex <file>; bibtex <file>; pdflatex <file>; pdflatex <file>)

\usepackage[T1]{fontenc}
\usepackage{lmodern}
\usepackage[margin=1in]{geometry}
\usepackage{microtype}
\usepackage{amsmath,amssymb,amsthm,mathtools}
\usepackage{booktabs,tabularx,array}
\usepackage{enumitem}
\usepackage[dvipsnames]{xcolor}
\usepackage[authoryear,round]{natbib}
\usepackage{xurl}  % lets long URLs break anywhere (loads url)

% Allow a line break (without an inserted hyphen) after each \_ in text, so long file names can wrap.
\makeatletter
\DeclareRobustCommand{\_}{\ifmmode\nfss@text{\textunderscore}\else\textunderscore\allowbreak\fi}
\makeatother
\usepackage[colorlinks=true,linkcolor=NavyBlue,citecolor=ForestGreen,urlcolor=RoyalBlue]{hyperref}
\usepackage[capitalise,noabbrev]{cleveref}

% ---- theorem-like environments (defined after cleveref so shared counters get the right names)
\theoremstyle{plain}
\newtheorem{theorem}{Theorem}
\newtheorem{proposition}[theorem]{Proposition}
\newtheorem{lemma}{Lemma}
\newtheorem{corollary}{Corollary}
\theoremstyle{definition}
\newtheorem{definition}{Definition}
\newtheorem*{assumptionU}{Assumption U}
\theoremstyle{remark}
\newtheorem{remark}{Remark}

% ---- status tags (robust, so they can appear in section titles and the table of contents)
\DeclareRobustCommand{\statustag}[2]{\textsf{\footnotesize\textcolor{#1}{[#2]}}}
\DeclareRobustCommand{\tagchecked}{\statustag{ForestGreen}{checked}}
\DeclareRobustCommand{\tagproved}{\statustag{RoyalBlue}{proved}}
\DeclareRobustCommand{\tagderived}{\statustag{RoyalBlue}{derived}}
\DeclareRobustCommand{\tagconj}{\statustag{BurntOrange}{conjecture}}
\DeclareRobustCommand{\tagopen}{\statustag{BrickRed}{open}}

% ---- research-problem headings (numbered P1, P2, ... by hand, listed in the table of contents)
\newcommand{\researchproblem}[2]{%
  \subsubsection*{#1.\ #2}\addcontentsline{toc}{subsubsection}{#1.\ #2}}

% ---- math and text macros
\newcommand{\E}{\mathbb{E}}
\newcommand{\R}{\mathbb{R}}
\newcommand{\N}{\mathbb{N}}
\newcommand{\Z}{\mathbb{Z}}
\newcommand{\Unif}{\mathcal{U}}
\DeclareMathOperator{\Var}{Var}
\DeclareMathOperator{\Corr}{Corr}
\DeclareMathOperator{\clip}{clip}
\DeclareMathOperator{\diag}{diag}
\DeclareMathOperator{\sgn}{sign}
\DeclareMathOperator{\fracpart}{frac}
\DeclareMathOperator{\rms}{rms}
\newcommand{\KL}{\mathrm{KL}}
\newcommand{\TV}{\mathrm{TV}}
\newcommand{\script}[1]{\texttt{sanity\_checks/\allowbreak#1}}
\newcolumntype{L}{>{\raggedright\arraybackslash}X}
\setlist{itemsep=2pt,topsep=4pt}

% ---- layout: avoid overfull lines from unbreakable math/citations; ragged-right bibliography
\setlength{\emergencystretch}{1.5em}
\renewcommand{\bibfont}{\raggedright}

\usepackage{flafter} % floats never appear before their place in the text
\usepackage{xltabular} % the status table may break across pages

\title{P1 Experiments:\\Testing the UAQ and Update-Tracking Theory}
\author{Experiments report for research problem P1 (draft)}
\date{18 November 2026 (split from the theory note on 12 November)}

\begin{document}
\maketitle

\paragraph{Scope.} This is the experiments report for research problem P1 in the companion report \texttt{research\_problems.tex}, and the companion to the theory note \texttt{p1\_uaq\_and\_update\_tracking.tex}, which states and proves the results tested here. Theorem, lemma, proposition, corollary, remark and assumption numbers refer to the theory note, as do QuRL's claims~(i) and~(ii) \citep{li2026qurl}. Float64 checks are in \texttt{sanity\_checks/}; GPU experiments in \texttt{experiments/}, all run on one H100.

\section*{In short}\addcontentsline{toc}{section}{In short}
\begin{itemize}
\item \textbf{Float64 checks (T)} agree with Theorems~1--2, Corollary~3 and Propositions~4--7.
\item \textbf{On vLLM's kernels (S1)}, Theorem~1 holds bit-for-bit: with INT8 the log-probabilities are identical at $s=2$ for 0.5B, 1.5B and 7B models, and at the paper's $s=1.5$ the mismatch does not shrink.
\item \textbf{In a training loop (S2)}, Theorem~2 holds bit-for-bit, in a QuRL-style loop and in VeRL; clipping rule~(c) is necessary.
\item \textbf{In RL (S3)}, QuRL's GSM8K result (Table~1) did not reproduce at the stated learning rate $10^{-5}$ (one run per format; most likely the critic's rate, as in VeRL's quickstart); at $10^{-6}$ the setup learns. The pre-registered RL test of ``UAQ $=$ matched optimizer settings'' is complete: 22 runs, R1--R6 for seeds 1--3 and R4, R5 for seeds 4--5. All three hypotheses are inconclusive on late accuracy, as the power analysis predicted; for H1 (R4 against R5, five seeds) the two R5 stalls and the seed-to-seed spread both stand in the way. On mismatch H2 is supported: UAQ changed the INT8 mismatch by about the same factor as the BF16 one (ratio of ratios 1.10, 90\% CI $[1.015, 1.20]$), so it does not reduce the quantization mismatch during RL either. Early stalls hit 4 of the 22 runs, all with INT8 rollouts. The INT8 rollout model follows the trainer.
\item \textbf{INT8 rollouts cost about 5 points here (S3, exploratory).} With BF16 rollouts the setup matches the published baselines (R1: 54.8--56.6\% late accuracy over three seeds; the paper's 55.35\%, VeRL's 56.7\%). QuRL's INT8 rollouts trail by 4.9 points averaged over UAQ (95\% CI $[2.3, 7.6]$), against the paper's 1.8, and every BF16 run ends above every INT8 run. The gap is 2--4 points at the first validation and 4.5--7.5 points from step 58 on, with lower training reward, higher entropy and more drift from the reference model. INT8 accuracy is scored on the INT8 engine.
\item \textbf{On a real QuRL run (E2, pre-registered)}, the theory's tracking and persistence results hold on the synced master weights in 87 of 88 checks.
  \begin{itemize}
  \item After 435 steps the rounding error still correlates 0.69 with its start, as Proposition~6 predicts.
  \item From the second update on, the rollout captures 97--99\% of each update. The one failure is the first update from the BF16 checkpoint, captured at 93\%; an exploratory simulation reproduces it from BF16 storage alone.
  \item In policy space the INT8 rollout follows the trainer with slope 0.93 over the run, where a frozen rollout gives 0. The registered $\pm10\%$ band fails in 2 of 11 checks.
  \item On stored checkpoints, Assumption~U holds at the rounding boundary in all 21 cells (7 projection types of the 0.5B, 1.5B and 7B models). Its bulk test fails narrowly in one cell (0.052 against 0.05), again from BF16 storage. Lemma~1's condition $\sigma/\Delta\ge1$ holds in every row.
  \end{itemize}
\item \textbf{Overall} (\cref{sec:assessment}), nothing in the experiments contradicts a theoretical result where its assumptions hold. E2's one registered weight-space failure, at the first update, is where BF16 storage breaks Assumption~U, by an exploratory analysis. The largest remaining gap is E3, what drives the mismatch rise, which is left to P4. The INT8 accuracy gap is what E3 and P6 would have to explain and close.
\end{itemize}

\tableofcontents

% =====================================================================================
\section{What is tested}\label{sec:tested}
The theory note answers five questions (its Section~2); the stages below test the answers. The algorithm under test is Algorithm~1 of the QuRL explainer (\texttt{qurl\_explained.tex}).
\begin{itemize}
\item \textbf{Q1. Does UAQ change the rollout model?} Theorem~1: no, for QuRL's W8A8 and any quantizer that is homogeneous on its groups. Tested in float64 (T) and on vLLM's kernels (S1).
\item \textbf{Q2. What does UAQ change in training?} Theorem~2 and Corollary~3: only optimizer settings, chiefly per-layer learning rates ($\times s$ on Q/K/V and $\div s$ on the norm gain under AdamW, with $\varepsilon$, weight decay and the clipping norm matched). Tested in float64 (T), bit-for-bit in a training loop and in VeRL (S2), and in distribution over seeds in full RL runs (S3, hypotheses H1--H3).
\item \textbf{Q3. Do re-quantized master weights freeze?} Propositions~3--5: no; updates are captured in expectation. Tested in float64 (T), by a first look in S3, and in full on a real run (E2).
\item \textbf{Q4. What is the real issue?} Propositions~6--7: the rounding error persists across steps, and error feedback bounds its cumulative effect. Tested in float64 (T) and on a real run (E2); E3 is left to P4.
\item \textbf{Q5. Do the answers hold on a real RL stack?} Stages S1--S3, E2 and E3 below.
\end{itemize}

\begin{center}\small
\begin{tabularx}{\linewidth}{@{}l l L >{\raggedright\arraybackslash}p{3.6cm}@{}}
\toprule
Stage & Question & What runs & Status\\
\midrule
T & Q1--Q4 (the maths) & float64 checks: one layer or a toy QuRL loop & done\\
S1 & Q1 on real kernels & step-0 identity test in vLLM, three model sizes & done\\
S2 & Q2, exactly & deterministic 20-step QuRL loop, UAQ vs.\ Theorem~2 settings at $s=2$ & done\\
S3 & Q5: H1--H3 & GSM8K PPO in VeRL: R1--R6 $\times$ 3 seeds ($+2$ for R4, R5); Table~1 reproduction & done (22 of 22)\\
E2, E3 & Q3--Q4 on real runs & tracking diagnostics (E2) and mismatch attribution (E3) & E2 done; E3 left to P4\\
\bottomrule
\end{tabularx}
\end{center}

% =====================================================================================
\section{Design}\label{sec:design}
All GPU stages run on one H100 NVL (94\,GB); T runs on a CPU.

\subsection{T: float64 checks}\label{sec:t}
\begin{itemize}
\item \script{uaq\_equivalence.py}: Theorem~1 for INT8 and FP8 at $s\in\{0.7,1.5,2,4\}$; Theorem~2 on a 300-step supervised toy (Adam and SGD).
\item \script{uaq\_edge\_cases.py}: the scope table of Remark~2.
\item \script{uaq\_pipeline\_equivalence.py}: Corollary~3 in a toy QuRL loop with INT4 rollouts, GRPO, TIS, ACR, AdamW and global-norm clipping (2 epochs $\times$ 4 minibatches, so ACR and clipping both bind).
\item \script{rtn\_tracking.py}: Propositions~4--5.
\item \script{tracking\_schemes.py}: Propositions~6--7 over 2{,}000 steps.
\end{itemize}

\subsection{S1: step-0 identity test on vLLM kernels}\label{sec:s1}
\begin{itemize}
\item \textbf{Models.} Qwen2.5-0.5B-Instruct \citep{qwen2024qwen25} (S3's model), Qwen2.5-Math-1.5B and Qwen2.5-Math-7B (the paper's size) \citep{yang2024qwen25math}, each with UAQ'd checkpoints at $s=2$ (exact in BF16) and $s=1.5$ (the paper's value; BF16 storage re-rounds about 4\% of INT8 codes).
\item \textbf{Rollout formats.} INT8 W8A8 on vLLM's \citep{kwon2023vllm} CUTLASS INT8 kernel (what QuRL would use on Hopper); FP8 per-tensor and per-channel as controls; BF16.
\item \textbf{Protocol.} Sample 512 GSM8K responses from the BF16 model; teacher-force every token under each configuration; compare per-token log-probabilities. Noise floor: the same checkpoint with shuffled batches.
\item \textbf{Metrics.} Identity: mean $|\Delta\log p|$ and share of identical tokens. Mismatch to the matched BF16 policy: mean $|\log\mu-\log\pi|$ (primary) and the sampled-token estimates $k_3$ and $k_1$ of $\KL(\pi\|\mu)$ on tokens sampled from $\pi$, with paired bootstrap CIs.
\item \textbf{Prediction.} $s=2$: identical. $s=1.5$: no systematic change. An $s$-fold error reduction would give ratios of 0.667 (mean $|\Delta\log p|$) and 0.444 (KL).
\end{itemize}

\subsection{S2: exact equivalence in a training loop}\label{sec:s2}
\begin{itemize}
\item \textbf{Why $s=2$.} Every rescaling is exact, so Theorem~2 predicts bit-identical runs.
\item \textbf{Loop.} Qwen2.5-0.5B-Instruct on GSM8K: 20 rollout batches of 8 prompts $\times$ 4 samples; GRPO with decoupled PPO, TIS ($C=2$) and ACR; 2 epochs $\times$ 2 minibatches (80 AdamW updates); $\lambda=0.01$, $\varepsilon=10^{-8}$, clipping at 0.1. Rollouts use INT8 W8A8 fake quantization in PyTorch with seeded sampling (vLLM cannot re-quantize INT8 online and is not repeatable at this size); training uses BF16 autocast with FP32 master weights and deterministic kernels.
\item \textbf{Two learning rates.} $\eta=10^{-6}$ (QuRL's regime; there weight decay changes no FP32 weight, since $\eta\lambda<2^{-25}$) and $10^{-5}$ (decay active).
\item \textbf{Arms.} R4 (UAQ, plain AdamW); R4 repeat; R5 (no UAQ, Theorem~2(b)--(c) settings); R5o (R5 with ordinary clipping); R5x (learning rates matched only); R3 (no UAQ, plain settings).
\item \textbf{Measure.} SHA-256 fingerprints at every step of the sampled tokens, $\mu$, $\pi_{\mathrm{old}}$ and the weights mapped to original coordinates.
\end{itemize}

\subsection{S3: QuRL's GSM8K PPO setup in VeRL}\label{sec:s3}
\begin{itemize}
\item \textbf{Configuration} (the paper's GSM8K setup, which is VeRL's quickstart, \citealp{sheng2025hybridflow}). Qwen2.5-0.5B-Instruct; PPO with a critic; batch 256, mini-batch 64; responses up to 512 tokens; 15 epochs $=$ 435 steps; greedy validation on the 1{,}319 test problems \citep{cobbe2021gsm8k} every 29 steps; KL penalty and entropy bonus 0.001, as in VeRL's documented baseline (56.7\% at $\eta=10^{-6}$). VeRL v0.9.1 with vLLM 0.24.
  \begin{itemize}
  \item The actor's gradient is clipped to global norm 1.0. The clip binds on 99.8--100\% of steps in every run (median pre-clip norm 2.5--4.7), so rule~(c) acts at every step.
  \item The critic starts from the original model, with the same optimizer, in every arm. UAQ changes only the starting checkpoint of the actor and of the frozen reference.
  \item The mismatch $k_3$ is VeRL's: an estimate of $\KL(\mu\|\pi)$ on the rollout tokens, which are sampled from $\mu$.
  \end{itemize}
\item \textbf{Added to VeRL.} INT8 W8A8 rollouts re-quantized at every sync (CUTLASS INT8 kernel); ACR as a policy loss; the UAQ checkpoint; an optimizer plug-in for R5 (per-group learning rate, $\varepsilon$ and decay, plus rule (c) clipping). All arms run the same code.
\item \textbf{Table~1 reproduction.} R1 (BF16) and R3 (QuRL INT8) at the paper's $\eta=10^{-5}$, one seed each. Targets: 55.35\% and 53.55\%.
\item \textbf{UAQ experiment.} At $\eta=10^{-6}$ (all the paper's UAQ results use $10^{-6}$) and $s=1.5$:
\begin{center}\small
\begin{tabular}{@{}lll@{}}
\toprule
Arm & Rollout & Change\\
\midrule
R1 & BF16 & ---\\
R2 & BF16 & UAQ\\
R3 & QuRL: INT8 W8A8, TIS ($C=2$), ACR & ---\\
R4 & QuRL & UAQ\\
R5 & QuRL & Theorem~2(b)--(c) settings instead of UAQ\\
R6 & QuRL & global learning rate $\times1.5$ (the paper's control)\\
\bottomrule
\end{tabular}
\end{center}
\item \textbf{Seeds.} 1--3 for every arm, plus 4--5 for R4 and R5 (added on 2026-11-11 for power, before any R5 result was seen): 22 runs of 4.2--5.5\,h, 104 H100-hours in all. Seed $k$ fixes the prompt order and the sampling seed, so arms are paired by seed.
\item \textbf{Gate.} The queue first repeats R5's exactness check on the H100 (three updates on one fixed batch at $s=2$) and starts only if it passes.
\item \textbf{Hypotheses.}
  \begin{itemize}
  \item \textbf{H1:} R4 and R5 are equal in distribution. Corollary~3 makes them identical in exact arithmetic; at $s=1.5$, that the rounding acts like a change of seed is a hypothesis (Remark~3).
  \item \textbf{H2:} UAQ's effect is not specific to quantization: (R2 $-$ R1) $\approx$ (R4 $-$ R3).
  \item \textbf{H3:} R6 differs from R4: a global learning-rate increase is the wrong control.
  \end{itemize}
\end{itemize}

\subsection{Pre-registered analysis of S3}\label{sec:prereg}
Fixed on 2026-11-11: H1 before any R5 result, H2 and H3 before any R1, R2 or R6 run. The scripts are frozen by SHA-256 and were tested on synthetic runs with known answers; any later change is recorded here, and results from both versions are reported.
\begin{itemize}
\item \textbf{Primary endpoint.} Late accuracy: mean greedy GSM8K test accuracy over the last five validations (steps 319--435), in points.
\item \textbf{Secondary endpoints} (supportive only). Final accuracy (step 435); reward AUC (mean training score over all steps); mismatch (mean $k_3$, compared as a ratio with bounds 0.80--1.25, with the entropy difference reported next to it); for H2 and H3 also early reward (mean score over steps 1--58, added after R4 seed~1's fast start had been seen).
\item \textbf{Margin.} $\delta=2.0$ points: the gap the paper calls small (``reduces this gap to 2\%''), and about the largest difference in its UAQ ablation.
\item \textbf{Decision rules} (paired by seed; t-intervals on $n-1$ degrees of freedom):
\begin{center}\small
\begin{tabular}{@{}lllll@{}}
\toprule
Hypothesis & Contrast per seed & Seeds & Supported if & Contradicted if\\
\midrule
H1 & R4 $-$ R5 & 1--5 & 90\% CI within $\pm\delta$ & 95\% CI excludes 0\\
H2 & (R4 $-$ R3) $-$ (R2 $-$ R1) & 1--3 & 90\% CI within $\pm\delta$ & 95\% CI excludes 0\\
H3 & R6 $-$ R4 & 1--3 & 95\% CI excludes 0 & 90\% CI within $\pm\delta$\\
\bottomrule
\end{tabular}
\end{center}
Otherwise inconclusive. If both conditions hold: ``equivalent within $\delta$, smaller difference detected'' (H1, H2) or ``supported, with a difference smaller than $\delta$'' (H3). For the mismatch, H2's contrast is the ratio of ratios (R4/R3)/(R2/R1): if UAQ helped the quantized rollout capture updates, R4/R3 would be below 1 and R2/R1 near 1.
\item \textbf{Also reported, without verdict.} UAQ's own effect R4 $-$ R3 (if it is within $\pm\delta$, H1 cannot distinguish ``R5 reproduces UAQ'' from ``UAQ does nothing here''), R2 $-$ R1, and R6 $-$ R3.
\item \textbf{Data rules.} Every completed run counts; a crashed run is re-run with the same seed; no verdict before all registered runs are complete.
\item \textbf{Power.} At zero true difference, H1's equivalence is reachable only if the paired SD is below about 1.2 points with 3 seeds, or 2.1 with 5. H2 needs a single-run spread below about 0.6 points, so it will likely be inconclusive. H3 detects only differences above about 2.5 paired SDs.
\item \textbf{Exploratory} (fixed after R5 seed~1's stall was seen at step 316, before any other R4 or R5 seed). A run counts as collapsed if its late accuracy is below 10\%; collapses are counted per arm, and H1--H3 are repeated without collapsed seeds. Each run's longest stretch without a correct training answer is reported.
\item \textbf{Seen before fixing.} H1: R3 seed~1 in full, one progress line of R4 seed~1 (step 44), nothing from R5. H2/H3: R3 and R4 seed~1 in full, R5 seed~1 through step 316; for seed~1, R4 $-$ R3 was $-1.2$ points in late accuracy and $+7.4$ in early reward.
\item \textbf{Code.} \texttt{experiments/s3/h1\_analysis.py}, run with \texttt{-{}-seeds 1 2 3 4 5}; SHA-256:\\ \texttt{\small 5ac6da470e152a7b1680921bdcb2c1281d09cc52b63ce015f144c9c05f368f6f}.\\ \texttt{experiments/s3/h23\_analysis.py}, which reuses \texttt{h1\_analysis.py}'s arithmetic; SHA-256:\\ \texttt{\small 5ed1a3b280b340550b009e0e54c5d95ee2e00a00f2bddd74b7fa612af1dfb3e4}.\\ Each report repeats the hashes.
\end{itemize}

\subsection{E2: tracking and persistence on real weights (pre-registered 2026-11-15)}\label{sec:e2}
Fixed before any E2 data were analysed. \emph{Seen beforehand:} the smoke test's dump check (syncs 1--4; it reports only that the dumps match the rollout's codes and the largest weight change, about $4\times10^{-6}$ per step), and a trial run of the static script on Qwen2.5-0.5B-Instruct, so the 0.5B checkpoint's results do not count as a fresh test.
\begin{itemize}
\item \textbf{What it tests.} Assumption~U and Propositions~3--7 of the theory note on the synced FP32 master weights of a real QuRL run, where T used float64 simulations. P3 and P7's error-feedback bound are exact and serve as consistency checks.
\item \textbf{Data.}
  \begin{itemize}
  \item \emph{Static:} the decoder linear weights of Qwen2.5-0.5B-Instruct, Qwen2.5-Math-1.5B and Qwen2.5-Math-7B, as stored (BF16).
  \item \emph{Dynamic:} one run of arm R3 with S3's own call (seed~1, $\eta=10^{-6}$, 435 steps), with a vLLM worker extension (\texttt{experiments/e2/e2\_dump.py}) that records the synced FP32 master weights. It saves the full state at syncs 1, 2, 3, 5, 9, 17, 33, 65, 66, 129, 257, 258, 435 and 436 (sync $k$ holds the weights after step $k-1$), and 16 fixed rows of every linear weight in layers 0, 6, 12, 18 and 23 at every sync.
  \item \emph{Smoke test:} the dumps are the weights the rollout quantized; layer~0's live INT8 codes equal the quantizer applied to the dump at every sync.
  \end{itemize}
\item \textbf{Quantities.} All in grid units, with the rollout quantizer's arithmetic: $\Delta=$ row absmax $\times\mathrm{fl}_{32}(1/127)$, $q=w/\Delta$, $c=\mathrm{round}(q)$, phase $\varphi=q-c$. Each row's absmax element is left out. Every sync uses its own grid, so the exact change $u$ of $q$ between syncs includes the grid's motion when a row's absmax moves. To first order that motion is $-(c+\varphi)\,\delta\Delta/\Delta$, which for large codes exceeds the weight's own update.
\item \textbf{Predictions and decision rules.}

\smallskip
{\footnotesize
\begin{tabularx}{\linewidth}{@{}l L L@{}}
\toprule
Claim & Measurement & Rule\\
\midrule
U1: U near the boundary & share of $|\varphi|\ge0.45$, divided by 0.10 & in $[0.95,1.05]$ for each projection type of each checkpoint, and at each full-dump sync\\
U2: U in the bulk, small codes & largest deviation of 40 bin densities from 1, $|c|\le15$ & $\le0.05$, as for U1\\
Lemma~1 & share of rows with $\sigma/\Delta<1$ & below 1\% (static)\\
P3: bounded lag & $\max|e_k-e_0|/\max(\Delta_k,\Delta_0)$ & $\le1$ (exact)\\
P4a: one-step capture & flip share $F$ over $P=\mathrm{mean}\min(1,|u|)$, $u=q_{k+1}-q_k$ & in $[0.90,1.10]$: with all elements for the pairs from sync 65 on, and for every pair without the elements that were exact ties at sync~1\\
P4b & observed against predicted flip rate in bins of $|u|$ (slope through the origin) & in $[0.90,1.10]$\\
P4c: captured in expectation & $\sum\Delta c\,u/\sum u^2$ & in $[0.95,1.05]$\\
P6a: stationary size & $\Var(\varphi_k)/\Var(\varphi_0)$ & in $[0.98,1.02]$ at every full-dump sync\\
P6b: persistence & $\Corr(\varphi_k,\varphi_0)$ in 20 bins of $a=\fracpart|q_k-q_0|$, against the bin mean of $1-6a(1-a)$ & largest deviation $\le0.05$ over bins with $\ge10^4$ elements\\
P6c & the same along $k$, from the row samples & largest deviation $\le0.05$ over syncs\\
P7a & error feedback's partial sums & $|\sum e_k|\le\Delta_K/2$ (exact)\\
P7b & stochastic rounding's partial sums & rms over $\sqrt{K/6}$ in $[0.90,1.10]$ at the last sync\\
P5: tracking in policy space & OLS slope of $\Delta\log\mu$ on $\Delta\log\pi$ over probe tokens & in $[0.90,1.10]$ for every cumulative pair (sync 1, $s'$) with $s'\ge9$ and for the pooled one-step pairs\\
P5$'$: not frozen & the same slope & $>0.5$\\
\bottomrule
\end{tabularx}\par}
\item \textbf{Policy-space measurement.} The probe set is 128 GSM8K training prompts (a fixed random subset), each with one response sampled at temperature~1 from $\theta_0$. For each full-dump sync, $\log\pi$ comes from an FP32 forward pass with the dumped weights. $\log\mu$ comes from the same model with the rollout's W8A8 quantizer applied as deterministic FP32 fake quantization: per-row INT8 weights and per-token INT8 activations, with embeddings, norms and the output head unquantized.
\item \textbf{Deviations expected in advance.}
  \begin{itemize}
  \item \emph{BF16 lattice.} Near a row's maximum, BF16 spacing is 0.5--1 grid step, so weights with large codes sit on a coarse lattice. U2 is therefore claimed only for $|c|\le15$.
  \item \emph{Exact ties.} About 0.2\% of the stored weights sit exactly on a rounding boundary, and any update flips about half of them. So $F/P$ exceeds 1 in the first pairs with all elements, and the rule for those pairs leaves the ties out. By sync 65 the ties have drifted apart, and the rule uses all elements.
  \item \emph{Moving grids.} $u$ includes the grid's motion. The fixed-grid prediction $\mathrm{mean}\min(1,|\delta|/\Delta)$ is reported next to it, so the size of that motion is visible.
  \end{itemize}
\item \textbf{Reading.} A failure of P4 or P6 where U1 holds would contradict Proposition~4 or~6 as stated. A failure of P5$'$ would support QuRL's claim~(i), that the rollout freezes. A failure of U limits where the propositions apply rather than refuting them.
\item \textbf{Also reported, without rules.} $\nu=\eta/\Delta$ per layer at 0.5B, 1.5B and 7B (the link to 7B); the correlation of phase and update; RTN's and error feedback's partial sums; and each one-step pair's slope.
\item \textbf{Code (frozen).} In \texttt{experiments/e2/}; each script has a self-test with known answers. On Gaussian weights all rules pass; on weights placed on the grid, U1 and P4 fail. A gradient step gives a policy-space slope of 1.03, and frozen rollout weights give 0.00. SHA-256:\\
\texttt{e2\_static.py} \texttt{\small 0db4ecece22b552caa49cecc83dbac474f8a929db7748d8bb7121af194cb40c4};\\
\texttt{e2\_dynamic.py} \texttt{\small 9ec9dca1f693962b3f7ffe1a2cdd3a2dc930c469de0ba857656451f7f2556fb3};\\
\texttt{e2\_policy.py} \texttt{\small af8f517022cda4b496af6425a4558480d6975750a2871a50a0d3695c7c2ec49c}.
\end{itemize}

\subsection{E3, and the link to 7B}\label{sec:link7b}
\begin{itemize}
\item \textbf{E3}, logged during training: per-layer $\|e_k\|/\|\theta_k\|$, weight and activation absmax, exact KL by position and by token probability. What drives the mismatch growth is research problem P4's question, so E3 is left to P4. Its one P1 part, the stationary size of the weight error, is P6a of E2 (\cref{sec:e2}).
\item \textbf{Link to 7B.} Early in Adam training each element moves by about $\eta$ per step, so $\nu_l\approx\eta/\Delta_l$, computable from the 7B checkpoint without training it. Comparing the distributions of $\nu_l$ shows whether a smaller model's runs are in the 7B model's tracking regime (Proposition~4 depends only on $|d|/\Delta$).
\item \textbf{Why not paper scale.} A 7B full fine-tune needs about 112\,GB for weights, gradients and Adam state alone; R3--R5 $\times$ 3 seeds would take 1{,}400--2{,}700 H100-hours. H1 at 7B is future work.
\end{itemize}

% =====================================================================================
\section{Results}\label{sec:results}

\subsection{T: float64 checks}
\begin{center}\small
\begin{tabularx}{\linewidth}{@{}l L@{}}
\toprule
Claim & Result\\
\midrule
Theorem~1 (INT8, FP8) & output change $\le2.9\times10^{-16}$ for $s\in\{0.7,1.5,2,4\}$; a per-channel (SmoothQuant, \citealp{xiao2023smoothquant}) vector does change the INT8 error, 5.2\% $\to$ 1.5\%\\
Theorem~2, supervised toy & Adam $1.6\times10^{-6}$ (limited by unadjusted $\varepsilon$), SGD $10^{-14}$; first-step $\nu$ gain 1.500 (Adam) vs.\ 2.250 (SGD)\\
Corollary~3, toy QuRL loop & matched settings: 100\% identical actions, weights within $3\times10^{-10}$; ordinary clipping: 59.7\% identical actions; plain settings: 50.1\%\\
Props.~4--5 & flip rate $=|d|/\Delta$ from $10^{-3}$ to 0.3; slope $\approx1$; aligned correlation 0.75 at $\epsilon=10^{-3}$, $n=2048$ (predicted 0.75)\\
Props.~6--7 & correlation within 0.007 of $1-6a(1-a)$; variance ratio 0.995--1.007 over 2{,}000 steps; rms partial sums after 2{,}000 steps: $292\Delta$ (RTN), $576\Delta$ (delta), $18\Delta$ (SR), $0.29\Delta$ (error feedback)\\
\bottomrule
\end{tabularx}
\end{center}

\subsection{S1: UAQ does not change the rollout model on real kernels}\label{sec:s1-results}
\textbf{At $s=2$} (UAQ vs.\ original, per-token log-probabilities):
\begin{center}\small
\begin{tabularx}{\linewidth}{@{}l L L l@{}}
\toprule
Format & 0.5B & 1.5B & 7B\\
\midrule
\textbf{INT8 W8A8} (QuRL) & \textbf{identical} & \textbf{identical} & \textbf{identical}\\
FP8 per-tensor / per-channel & identical & run-to-run variation (an $s=2$ run reproduced an original run bit-for-bit) & identical\\
BF16 & 0.36\% of tokens differ (engine nondeterminism) & identical & identical\\
\bottomrule
\end{tabularx}
\end{center}

\textbf{At $s=1.5$} (mismatch with UAQ $\div$ mismatch without, 95\% CI; an $s$-fold reduction would be 0.667 and 0.444):
\begin{center}\small
\begin{tabular}{@{}llll@{}}
\toprule
Format & Model & mean $|\Delta\log p|$ ratio & KL ($k_3$) ratio\\
\midrule
\textbf{INT8 W8A8} & 0.5B & \textbf{0.994} [0.984, 1.004] & \textbf{0.978} [0.951, 1.007]\\
\textbf{INT8 W8A8} & 1.5B & \textbf{0.996} [0.988, 1.005] & \textbf{0.994} [0.966, 1.022]\\
\textbf{INT8 W8A8} & 7B & \textbf{1.013} [1.004, 1.021] & \textbf{1.022} [0.995, 1.049]\\
FP8 per-tensor & 0.5B / 1.5B / 7B & 1.004 / 0.976 / 0.946 & 1.013 / 0.923 / 0.862\\
FP8 per-channel & 0.5B / 1.5B / 7B & 1.007 / 0.968 / 0.943 & 1.013 / 0.888 / 0.964\\
\bottomrule
\end{tabular}
\end{center}
\begin{itemize}
\item \textbf{Theorem~1 holds bit-for-bit at $s=2$} for QuRL's INT8 W8A8 at all three sizes.
\item \textbf{At the paper's $s=1.5$, the INT8 mismatch does not shrink at any size}, 7B included. The claimed error-reduction mechanism is absent at step 0 on the paper's kernel.
\item \textbf{FP8} shows small reductions at 1.5B and 7B, never near $s$-fold. They are three BF16 re-rounding draws, two favourable; per-token errors before and after UAQ correlate only 0.43--0.64. A null distribution over re-rounding draws would settle whether any of it is systematic.
\item \textbf{Side findings.} vLLM's default engine is neither batch-invariant nor run-to-run deterministic at 0.5B and 1.5B (0.22--0.31\% of tokens change under reshuffling, BF16 included); INT8 W8A8 was identical in every same-order comparison. The INT8--FP8 KL ratio grows with size (1.0--1.1$\times$ at 0.5B, 1.4--1.5$\times$ at 1.5B, 2.8--3.0$\times$ at 7B), consistent with activation outliers \citep{dettmers2022llmint8,xiao2023smoothquant} but not yet attributed. The mismatch is heavy-tailed: the top 1\% of tokens carry 38--48\% of the KL. Sampled-token KL estimates are fragile (rare tokens with $\mu/\pi$ up to 17--37), so mean $|\Delta\log p|$ is the primary metric.
\end{itemize}

\subsection{S2: Theorem~2 holds bit-for-bit}\label{sec:s2-results}
\begin{center}\small
\begin{tabularx}{\linewidth}{@{}l L l l@{}}
\toprule
Arm & Differs from R4 in & Bit-identical steps ($10^{-6}$; $10^{-5}$) & Max.\ weight diff.\ (20 steps)\\
\midrule
R4 repeat & nothing & 20/20; 20/20 & 0; 0\\
\textbf{R5} & coordinates only & \textbf{20/20; 20/20} & \textbf{0; 0}\\
R5o & clipping norm in original coordinates & 0/20; 0/20 & 1.6e-4; 1.1e-3\\
R5x & learning rates only matched & 0/20; 0/20 & 1.7e-4; 1.3e-3\\
R3 & plain settings, no UAQ & 0/20; 0/20 & 1.4e-4; 1.1e-3\\
\bottomrule
\end{tabularx}
\end{center}
\begin{itemize}
\item \textbf{R5 matched R4 in every fingerprint at every step}, and its saved weights equal R4's mapped weights element by element. The test covers INT8 W8A8 rollouts, BF16 autocast with FP32 master weights, GRPO, TIS, ACR, AdamW and global-norm clipping (binding on every update with a nonzero gradient), and weight decay at $10^{-5}$.
\item \textbf{Clipping rule (c) is necessary.} On the first update the gradient norm is 1.131 in UAQ coordinates and 0.874 in original ones, so ordinary clipping (R5o) diverged at step~1, as early as R3.
\item \textbf{One differing bit decorrelates the runs}, so at $s=1.5$, where $W/s$ is rounded, H1 must be tested in distribution over seeds.
\item \textbf{In VeRL}, the R5 plug-in passed the same test inside VeRL's FSDP2 actor on the H100 (three updates on one fixed batch at $s=2$): R5 and an R4 repeat were bit-identical to R4 in all 494M parameters; R5o differed from the first update.
\end{itemize}

\subsection{S3: QuRL's Table~1 is not reproduced at \texorpdfstring{$\eta=10^{-5}$}{eta = 1e-5}}\label{sec:t1-results}
\begin{center}\small
\begin{tabular}{@{}llll@{}}
\toprule
Run & GSM8K test accuracy at step 435 & Best of 15 validations & Paper (Table~1)\\
\midrule
R1, BF16 & 0.0\% & 0.0\% & 55.35\%\\
R3, QuRL INT8 & 2.9\% & 3.4\% (step 58) & 53.55\%\\
\bottomrule
\end{tabular}
\end{center}
\begin{itemize}
\item \textbf{Neither run learned.} R1's responses shrank to 2--3 tokens, with a KL from the reference model of about 5 nats per token. From step 290, R3's responses all ran to the 512-token cap at an entropy of 0.001.
\item \textbf{The learning rate is the cause.} The 10-step pilots of the same configuration at $10^{-6}$ show it: after 10 steps the KL from the reference model is $2.1\times10^{-3}$ at $10^{-6}$ against 0.15 at $10^{-5}$; the PPO clip fraction is 10--20$\times$ higher and the entropy rises from 0.50 to 0.80; correct answers in steps 5--10 number 6 against none.
\item \textbf{Mechanism.} At step 1 the critic's new value head predicts $-20$ to $+33$ for returns in $[0,1]$, so early advantages mostly reflect its error. At $10^{-5}$ the actor follows them ten times as fast and leaves the region where correct answers occur; afterwards only the KL penalty acts.
\item \textbf{Two explanations fit.} (1) The paper's actor ran at $10^{-6}$ and $10^{-5}$ is the critic's rate, as in VeRL's baseline (56.7\%, within 1.4 points of the paper's 55.35\%). (2) The paper used an unstated stabilizer such as warmup. Its text favors (2), its numbers fit (1); the paper does not let us decide.
\item \textbf{Side result.} A near-deterministic policy hides the mismatch: in R3 at $10^{-5}$, $k_3$ tracked the entropy (Spearman 0.94); in a normal run at $10^{-6}$ it does not (0.13).
\end{itemize}

\subsection{S3: the UAQ experiment at \texorpdfstring{$\eta=10^{-6}$}{eta = 1e-6}}\label{sec:uaq-results}
\textbf{Status (2026-11-15).} All 22 runs are complete: R1--R6 for seeds 1--3, and R4 and R5 for seeds 4--5, which only H1 uses. Every registered verdict is in. The gate passed on the H100: R5 and an R4 repeat were bit-identical to R4 in all 494M parameters; R5o differed in 99.98\% of $W$. The table lists every run.
\begin{table}[tbp]
\centering\footnotesize
\begin{tabularx}{\linewidth}{@{}l c c c c c c c L@{}}
\toprule
Run & Seed & Late & Final & Early rew. & Rew.\ AUC & $k_3$ & Entropy & Longest stretch without a correct answer\\
\midrule
R1, BF16 & 1 & 55.9 & 56.3 & 29.6 & 63.9 & 0.0006 & 0.12 & none\\
R2, BF16 + UAQ & 1 & 55.7 & 56.1 & 37.0 & 64.5 & 0.0006 & 0.11 & none\\
R3, QuRL INT8 & 1 & 50.2 & 50.1 & 34.4 & 56.0 & 0.055 & 0.26 & none\\
R4, QuRL + UAQ & 1 & 49.0 & 52.2 & 41.8 & 57.6 & 0.058 & 0.28 & none\\
R5, Thm.~2 settings & 1 & 11.8 & 46.6 & 0.1 & 2.9 & 0.016 & 0.44 & 370 steps (4--373)\\
R6, QuRL, lr $\times1.5$ & 1 & 48.7 & 49.4 & 41.3 & 55.6 & 0.037 & 0.24 & none\\
\midrule
R1 & 2 & 54.8 & 54.0 & 36.8 & 65.0 & 0.0007 & 0.13 & 1 step (step 2)\\
R2 & 2 & 54.6 & 52.8 & 36.4 & 64.3 & 0.0007 & 0.12 & none\\
R3 & 2 & 51.4 & 49.6 & 43.7 & 57.4 & 0.055 & 0.30 & none\\
R4 & 2 & 49.0 & 50.2 & 40.8 & 55.2 & 0.063 & 0.34 & none\\
R5 & 2 & 49.9 & 50.4 & 36.3 & 56.1 & 0.063 & 0.30 & 1 step (steps 5 and 10)\\
R6 & 2 & 51.2 & 50.0 & 31.1 & 55.1 & 0.041 & 0.35 & 5 steps (9--13)\\
\midrule
R1 & 3 & 56.6 & 57.5 & 37.1 & 64.9 & 0.0007 & 0.13 & 1 step (step 5)\\
R2 & 3 & 53.9 & 55.6 & 37.5 & 64.3 & 0.0006 & 0.11 & none\\
R3 & 3 & 51.0 & 51.9 & 2.9 & 44.1 & 0.068 & 0.28 & 21 steps (2--22)\\
R4 & 3 & 51.2 & 51.2 & 34.1 & 57.2 & 0.062 & 0.33 & 3 steps (4--6)\\
R5 & 3 & 44.5 & 48.4 & 0.1 & 43.8 & 0.040 & 0.29 & 72 steps (9--80)\\
R6 & 3 & 47.8 & 47.0 & 0.0 & 30.0 & 0.020 & 0.32 & 156 steps (25--180)\\
\midrule
R4 & 4 & 45.9 & 48.3 & 25.1 & 51.4 & 0.056 & 0.34 & 1 step (step 4)\\
R5 & 4 & 50.8 & 49.4 & 34.4 & 53.9 & 0.070 & 0.31 & 6 steps (4--9)\\
\midrule
R4 & 5 & 53.0 & 54.1 & 33.6 & 58.0 & 0.050 & 0.24 & none\\
R5 & 5 & 48.9 & 50.0 & 31.1 & 54.3 & 0.066 & 0.36 & 1 step (step 9)\\
\bottomrule
\end{tabularx}

\smallskip\raggedright
Late and final: accuracy in \%; early reward and reward AUC: mean training score in \%; $k_3$ and entropy: means over all steps. For reference: the paper's BF16 and QuRL INT8 results are 55.35\% and 53.55\%, VeRL's BF16 baseline 56.7\%. INT8 arms (R3--R6) are scored on the INT8 engine.
\end{table}
\begin{itemize}
\item \textbf{The setup learns at $10^{-6}$.} R3 reached 41.4\% and 46.4\% at step 29 in seeds 1 and 2 (VeRL's BF16 log: 45.8\%), and its final accuracy is 49.6--51.9\% over three seeds, 1.7--4.0 points below the paper's INT8 result. This confirms that the $10^{-5}$ failure came from the learning rate.
\item \textbf{BF16 matches the published baselines} (R1 and R2, seeds 1--3). R1's late accuracy is 54.8--56.6\% (mean 55.8) and its final accuracy 54.0--57.5\%, against the paper's 55.35\% and VeRL's 56.7\%; R2's late accuracy is 53.9--55.7\% (mean 54.8). No BF16 run stalled: the training score first reached 0.3 at steps 18--26. Their mismatch, a mean $k_3$ of $5.8$--$7.3\times10^{-4}$, is the BF16 floor between vLLM and the FSDP trainer.
\item \textbf{QuRL's INT8 rollouts cost about 5 points here} (exploratory; R1 against R3 was the planned check of the INT8 path).
  \begin{itemize}
  \item \textbf{Size.} R1 $-$ R3 in late accuracy is $+5.7$, $+3.4$ and $+5.6$ points per seed (mean $+4.9$, 95\% CI $[1.7, 8.1]$). Averaged over UAQ, that is R1 and R2 against R3 and R4 (the rollout format's main effect in H2's $2\times2$ design), it is $+4.9$ points (95\% CI $[2.3, 7.6]$), and $+4.5$ in final accuracy ($[2.4, 6.7]$).
  \item \textbf{Complete separation.} All six BF16 runs end above all sixteen INT8 runs, seeds 4--5 included, in late accuracy (53.9--56.6\% against at most 53.0\%). By chance alone that ordering has probability $1/\binom{22}{6}\approx1.3\times10^{-5}$.
  \item \textbf{Against the paper.} Its gap is 1.8 points. The main effect's 95\% CI excludes it; R1 $-$ R3's just includes it. If the paper reports best validations, the gap by best validation here is 3.7 points (R1 56.2--57.5\%, R3 52.5--53.8\%).
  \item \textbf{When it opens.} The twelve INT8 runs without an early stall start lower, though the ranges overlap: at step 29, 33.9--47.4\% against 43.2--47.6\% with BF16. On average the gap there is 2.1 points over seeds 1--3 and 4.0 with seeds 4--5. It widens to 5--7.5 points over steps 58--145 and stays at 4.5--6.7 points afterwards. Over steps 87--145, every BF16 validation is above every validation of those twelve runs (52.4--54.9\% against 39.3--50.9\%).
  \item \textbf{How the runs differ} (BF16 against those twelve INT8 runs; the same in every seed):
  \begin{itemize}
  \item training score: no gap before step 30; it opens over steps 30--116 (3.3, then 7.1 points between window medians) and stays at 9.4--10.9 points from step 117 on, ending at 72--73\% against 56--65\%;
  \item entropy: INT8 entropy bottoms out near 0.21 over steps 59--116 and rises again to about 0.3, while BF16 entropy keeps falling, to 0.05--0.07 at the end (INT8: 0.17--0.43);
  \item KL from the reference model at the end: 0.20--0.30 against 0.48--0.84;
  \item not response length: seed~1's BF16 runs kept longer answers, seed~2's did not (over steps 200--435, 166--279 tokens against 138--186).
  \end{itemize}
  \item \textbf{Updates follow the INT8 samples.} The PPO-KL diagnostic, the mean of $\log\pi_{\mathrm{old}}-\log\pi_\theta$ over the rollout's tokens, is positive with BF16 in every window median after step 10 (about $+10^{-4}$). In the INT8 runs all 60 window medians through step 174 are negative, down to $-2.3\times10^{-3}$. So on the INT8 sampler's tokens, the updated trainer puts more probability than $\pi_{\mathrm{old}}$ did. What drives this is not attributed.
  \item \textbf{Caveat: INT8 accuracy is scored on the INT8 engine.} VeRL validates greedily with the rollout engine, so an INT8 arm's accuracy is the quantized actor's. The runs save no checkpoints, so the trained weights cannot be scored in BF16 afterwards.
  \begin{itemize}
  \item \textbf{A fixed scoring penalty is not ruled out.} The INT8 runs are already 2--4 points lower at step 29, so part of the gap could be a fixed cost of greedy decoding with the INT8 policy.
  \item \textbf{But the runs also learn differently.} The gap grows to 5--7.5 points by step 145. The training reward, sampled at temperature~1 from each run's own rollout policy, shows no gap before step 30 and about 10 points from step 117 on.
  \item \textbf{To separate the two,} score the final actors in both formats.
  \end{itemize}
  \end{itemize}
\item \textbf{R5 seed~1 stalled for 370 steps, then learned.} Its few correct answers vanished within three steps while the entropy rose from 0.54 to 0.93 by step 10, the same early drift as the $10^{-5}$ runs. With no reward left, the KL penalty pulled the policy back toward the reference model (KL 1.1 at step 119, 0.03--0.05 by steps 300--390); correct answers returned at step 374, and validation went from 0.0\% (through step 377) to 46.6\% at step 435.
\item \textbf{No implementation fault.} The log shows the intended settings (Q/K/V: lr $1.5\times10^{-6}$, $\varepsilon$ $6.7\times10^{-9}$, decay 0.0067; norm gains: lr $6.7\times10^{-7}$, $\varepsilon$ $1.5\times10^{-8}$, decay 0.015; rest: $10^{-6}$; clip on the UAQ-coordinate norm), no non-finite gradients, and a constant scheduler that keeps each group's rate. Runs differ randomly from step 1 (R3 and R5 sampled different responses from identical weights), so a stall is a possible outcome of this configuration, not evidence that R5 differs from R4.
\item \textbf{R3 stalls too.} R3 seed~3 had no correct training answer at steps 2--22, the same pattern as R5 seed~1: correct answers vanish within a few steps, the entropy rises (0.52 $\to$ 0.85 by step 25) and the policy drifts from the reference model.
  \begin{itemize}
  \item \textbf{The stall.} Over steps 30--125 the median response was 7 tokens long and the KL from the reference model was above 1 at every step but one (peak 3.0 at step 93). Validation stayed at 5.5--8.3\% through step 116.
  \item \textbf{The recovery} came after about 125 steps instead of 370. Between steps 125 and 135 the policy returned to long answers (KL from the reference 2.1 $\to$ 0.10, training score 8\% $\to$ 42\%), and it validated at 49.2--52.6\% from step 174 on.
  \item \textbf{What it shows.} R3 runs VeRL's ordinary optimizer settings on the original checkpoint, so early stalls belong to this configuration at $10^{-6}$, not to R5's settings.
  \end{itemize}
\item \textbf{R5 seed~2 learned normally.} Correct training answers were missing only at steps 5 and 10; validation was 44.7\% at step 29 and 47.5--50.4\% afterwards. The log shows the same R5 settings as seed~1.
\item \textbf{R5 seed~3 stalled for about 80 steps, quietly.} It had 6, 3, 4 and 1 correct training answers at steps 1--4, then none at steps 5--7, 9--80 and 82. Validation was 0.0\% at steps 29 and 58, 25.5\% at step 87, and 42.3--48.4\% from step 116 on, apart from one dip. The log shows the same R5 settings as seeds 1 and 2.
  \begin{itemize}
  \item \textbf{Quiet.} Over steps 1--87 the policy stayed near the reference model (KL 0.006--0.047, window medians). Its entropy fell from 0.55 to 0.38, its responses stayed long (median 243--304 tokens), and the mismatch stayed low ($k_3$ $1.7$--$2.6\times10^{-2}$). R5 seed~1 and R3 seed~3 instead drifted far from the reference (KL about 1 or more) with rising entropy.
  \item \textbf{One late dip.} Validation fell to 33.6\% at step 319, between 45.1\% (step 290) and 46.3\% (step 348), while the training score held at 0.51--0.61. The dip came during the run's highest mismatch ($k_3$ 0.06--0.13 per step over steps 315--331). It alone moves seed~3's R4~$-$~R5 in late accuracy from 3.7 points (over the other four late validations) to 6.7.
  \end{itemize}
\item \textbf{R4 seed~3 wobbled but did not stall.} It had no correct answer at steps 4--6, the highest early entropy of any run in seeds 1--3 (0.96 at step 11; R4 seed~5 later reached 1.02), and a mismatch whose 29-step median reached 0.128 (steps 26--54). It recovered at once: 15 correct answers at step 10, a training score of 0.3 from step 22, and validation of 42.4--45.6\% through step 116 and 46.8--53.3\% afterwards.
\item \textbf{R6 seed~3 stalled for about 180 steps, quietly at first.} It had one correct training answer at step~1, then none at steps 2--23 and 25--180 (nor at 182--184, 186 and 188). Validation was 0.0\% through step 174, 38.0\% at step 203 and 45.6--48.7\% afterwards; its late accuracy is 47.8\%.
  \begin{itemize}
  \item \textbf{Quiet first, like R5 seed~3.} Through step 145 it stayed near the reference model (KL 0.013--0.095, window medians), with entropy 0.31--0.48.
  \item \textbf{Then a drift.} The KL rose to 0.53--0.57 over steps 146--203, before the recovery. Answers shortened from 317 to about 155 tokens.
  \item \textbf{Low mismatch throughout:} $k_3$ window medians stayed at most $2.7\times10^{-2}$.
  \end{itemize}
\item \textbf{Early stalls hit 4 of 22 runs, all with INT8 rollouts.} R5 seed~1, R3 seed~3, R5 seed~3 and R6 seed~3 validated at 0.0\%, 5.5\%, 0.0\% and 0.0\% at step 29; the other eighteen at 33.9--47.6\%. Seeds 4--5 had none.
  \begin{itemize}
  \item \textbf{By arm:} R3 1 of 3, R4 0 of 5, R5 2 of 5, R6 1 of 3; R1 and R2 none.
  \item \textbf{BF16 against INT8: suggestive only.} If the four stalls fell at random among the 22 runs, all four would land in the 16 INT8 runs with probability 0.25 (one-sided Fisher).
  \item \textbf{R4 against R5:} if both stalled equally often, both of their two stalls would land in R5 with probability 0.22.
  \item \textbf{Seed~3.} With the same prompt order, three of the four INT8 arms stalled, while both BF16 runs learned at once (early reward: R3 2.9, R4 34.1, R5 0.1, R6 0.0; R1 37.1, R2 37.5).
  \item \textbf{Same start, different drifts.} In each stall the few correct answers of the first steps vanished within ten steps, and none returned for a long stretch. With no task reward in that stretch, nothing task-driven steers the policy. R5 seed~1 and R3 seed~3 drifted far from the reference model with rising entropy; R5 seed~3 and R6 seed~3 stayed near it at first.
  \item \textbf{No collapse.} All four recovered. The lowest late accuracy is R5 seed~1's 11.8\%, above the 10\% threshold.
  \end{itemize}
\item \textbf{R4's faster start did not recur.} In seed~1, R4 had 49 correct answers per batch at step 10 against R3's 14 (early reward $+7.4$ points); in seed~2, 23 against 39 ($-3.0$); in seed~3, 15 against 0 ($+31.2$, as R3 stalled). Early reward varies widely between seeds of one arm (R3: 34.4, 43.7 and 2.9), so that endpoint is noisy.
\item \textbf{Side hypothesis, not supported.} The proposal was that the $1.5\times$ Q/K/V learning rate that R4 and R5 share (and R6 exceeds) makes the early phase more volatile: faster learning, and more stalls. R4 started slower than R3 in seed~2. R6, with every learning rate $\times1.5$, started fast in seed~1 (training score 0.3 from step 13), slower in seed~2 (step 25), and stalled in seed~3. Stalls hit R4 and R5 in 2 of 10 runs, R6 in 1 of 3 and R3 in 1 of 3: no higher rate with the larger Q/K/V rate.
\item \textbf{H1's registered verdict: inconclusive} (\texttt{h1\_analysis.py -{}-seeds 1 2 3 4 5}, frozen; 2026-11-15). R4 $-$ R5 over five pairs, $t$ on 4 degrees of freedom:

\smallskip
{\footnotesize
\begin{tabularx}{\linewidth}{@{}l L l l@{}}
\toprule
Endpoint & Per seed (1--5) & Mean, 90\% CI & Outcome\\
\midrule
Late accuracy (primary) & $+37.2$, $-0.9$, $+6.7$, $-4.9$, $+4.1$ & $+8.5$, $[-7.5, 24.4]$ & \textbf{inconclusive}\\
Final accuracy & $+5.5$, $-0.2$, $+2.8$, $-1.1$, $+4.0$ & $+2.2$, $[-0.5, 4.9]$ & inconclusive\\
Reward AUC & $+54.7$, $-1.0$, $+13.4$, $-2.5$, $+3.6$ & $+13.7$, $[-9.0, 36.3]$ & inconclusive\\
Mismatch ratio & 3.654, 0.997, 1.551, 0.807, 0.764 & 1.28, $[0.69, 2.38]$ & inconclusive\\
\bottomrule
\end{tabularx}\par}
  \begin{itemize}
  \item \textbf{Why.} Both large positive values come from R5's stalls (seeds 1 and 3). In the three seeds where neither arm stalled (2, 4 and 5), R4 $-$ R5 is $-0.9$, $-4.9$ and $+4.1$ points (exploratory), with no consistent sign.
  \item \textbf{Equivalence was out of reach even without stalls.} R4's late accuracy alone varied with an SD of 2.7 points over five seeds (45.9--53.0\%), so the paired SD below about 2.1 points that equivalence required (\cref{sec:prereg}) was unlikely.
  \item \textbf{``Contradicted'' did not happen.} It needed R4 $-$ R5 of about 31--35 points in both seeds 4 and 5; they gave $-4.9$ and $+4.1$.
  \item \textbf{Exploratory repeat.} No run collapsed (R5 seed~1's 11.8\% is above 10\%), so the repeat without collapsed seeds equals the registered test.
  \item \textbf{Mismatch.} R5's stalls had low mismatch, which pushes the ratio above 1. In seeds 4 and 5, R5's mismatch was the higher one (ratios 0.81 and 0.76).
  \end{itemize}
\item \textbf{The INT8 mismatch rises, then partly recedes.} $k_3$ starts at $1.02$--$1.16\times10^{-2}$ in all sixteen INT8 runs. In the ten H1 runs without a stall it reaches window medians of $5.5$--$8.4\times10^{-2}$ over steps 60--175, 4.7--7.8 times its start, and it ends (median of steps 400--435) at $3.6$--$6.3\times10^{-2}$ in every H1 run but R5 seed~1. That is about 40--100$\times$ the BF16 arms' $k_3$ over the same steps ($6.1$--$8.4\times10^{-4}$). R6 rose less: $4.1$ and $4.9\times10^{-2}$ over steps 60--175 in seeds 1 and 2, and $2.8$--$3.1\times10^{-2}$ at the end in all three seeds.
  \begin{itemize}
  \item \textbf{Not explained by entropy or drift.} In the twelve INT8 runs without a stall it tracks neither entropy (Spearman $-0.56$ to 0.13) nor the KL from the reference model ($-0.39$ to 0.38), so attribution needs E3.
  \item \textbf{The four stalls differ.} R3 seed~3's had the highest mismatch of any run: a 29-step median of 0.195 (steps 40--68), against at most 0.128 elsewhere (R4 seed~3's wobble), with TIS truncating a median 2.1\% of tokens per step (at most 3.4\%). The other three had the lowest: window medians at most $3.1\times10^{-2}$ (R5 seed~1), $2.6\times10^{-2}$ (R5 seed~3) and $2.7\times10^{-2}$ (R6 seed~3).
  \item \textbf{TIS and ACR.} In the other fifteen INT8 runs TIS truncated 0.19--2.07\% of tokens per step. The largest ratio in a run ranged from 153 (R5 seed~1) to 26{,}862 (R6 seed~2). ACR changed the clipping decision for 0.003--0.011\% of tokens.
  \item \textbf{Validation dips (exploratory).} How far a validation falls below its two neighbours correlates weakly with the mismatch over the 10 steps before it: Spearman $+0.20$ over 171 validations after learning in the INT8 runs ($+0.01$ to $+0.73$ per run).
  \begin{itemize}
  \item \textbf{Not specific to quantization.} The same weak correlation appears in the BF16 runs ($+0.18$ over 78 validations), where $k_3$ is only the numerical floor.
  \item \textbf{The three largest dips came at very different mismatch levels:} R5 seed~3's 12.1 points at step 319 at $1.5\times$ its run's median mismatch, R6 seed~1's 9.5 points at step 145 at $0.78\times$, and R5 seed~4's 8.4 points at step 232 at $1.01\times$.
  \item \textbf{INT8 curves are bumpier.} 7 of the 15 INT8 runs with validations after learning have a dip above 3 points, against none of the 6 BF16 runs (largest 2.5).
  \end{itemize}
  \end{itemize}
\item \textbf{UAQ's own effect, R4 $-$ R3, is small} (seeds 1--3, complete; reported without verdict, \cref{sec:prereg}).
  \begin{itemize}
  \item \textbf{Mismatch.} R4's mean $k_3$ is $1.036\times$ R3's (per seed 1.063, 1.134 and 0.923; 90\% CI $[0.867, 1.238]$). The CI lies inside the bounds 0.80--1.25, so UAQ changed the training mismatch by less than the margin, as S1 found at step 0 for $s=1.5$. H2's registered mismatch test confirms this against the BF16 arms (below).
  \item \textbf{Accuracy.} Late $-1.1$ points (90\% CI $[-3.3, 1.1]$), final $+0.7$ ($[-1.6, 3.0]$), reward AUC $+4.2$ ($[-9.4, 17.7]$). None of these CIs lies within $\pm\delta$ or excludes 0.
  \item \textbf{Entropy.} 0.017, 0.037 and 0.045 higher with UAQ.
  \item \textbf{Consequence for H1} (the caveat in \cref{sec:prereg}). With UAQ barely moving the mismatch, H1's mismatch endpoint cannot tell ``R5 reproduces UAQ'' from ``UAQ does little here''.
  \end{itemize}
\item \textbf{H2 and H3: the registered verdicts} (seeds 1--3, complete; \texttt{h23\_analysis.py}, frozen). No run collapsed, so the exploratory repeats without collapsed seeds equal the registered tests.

\smallskip
{\footnotesize
\begin{tabularx}{\linewidth}{@{}l L L@{}}
\toprule
Endpoint & H2: (R4 $-$ R3) $-$ (R2 $-$ R1) & H3: R6 $-$ R4\\
\midrule
Late accuracy (primary) & $-0.1$, $[-4.5, 4.3]$: \textbf{inconclusive} & $-0.5$, $[-5.2, 4.2]$: \textbf{inconclusive}\\
Final accuracy & $+1.7$, 95\% CI $[0.4, 3.0]$: difference detected & $-2.4$, $[-5.7, 1.0]$: inconclusive\\
Reward AUC & $+4.4$, $[-9.5, 18.2]$: inconclusive & $-9.8$, $[-35.3, 15.8]$: inconclusive\\
Early reward & $+9.4$, $[-21.9, 40.7]$: inconclusive & $-14.7$, $[-44.0, 14.5]$: inconclusive\\
Mismatch ratio & 1.10, $[1.015, 1.20]$: \textbf{supported} & 0.51, $[0.26, 1.00]$: inconclusive\\
\bottomrule
\end{tabularx}

\smallskip
Mean over seeds, in points (accuracy, reward) or as a ratio ($k_3$), with the 90\% CI unless a 95\% CI decided the outcome.\par}
  \begin{itemize}
  \item \textbf{H2 is inconclusive on its primary endpoint,} as the power analysis predicted (\cref{sec:prereg}). Equivalence needed an interaction SD below about 1.2 points; it was 2.6 (per seed $-1.0$, $-2.2$ and $+2.8$).
  \item \textbf{On mismatch, H2 is supported.} UAQ changed the INT8 mismatch by about the same factor as the BF16 one (R4/R3 1.04, R2/R1 0.94). If UAQ cut the INT8 error $s$-fold, the ratio of ratios would be about 0.44; its 90\% CI starts at 1.015.
  \item \textbf{On final accuracy, a small interaction was detected:} $+2.2$, $+1.7$ and $+1.1$ points per seed, with a mean of $+1.7$, inside $\pm\delta$. UAQ did slightly better on final accuracy with INT8 rollouts (R4 $-$ R3 $+0.7$) than with BF16 (R2 $-$ R1 $-1.0$). This does not change the verdict: it is one of four secondary endpoints, which are supportive only, and the primary endpoint, the mean of five validations, shows $-0.1$.
  \item \textbf{H3 is inconclusive on every endpoint.} R6 $-$ R4 in late accuracy is $-0.3$, $+2.2$ and $-3.4$ points per seed. The large AUC and early-reward differences come from R6 seed~3's stall. R6's mismatch was lower than R4's in every seed, but the ratio's CI reaches 1.00.
  \item \textbf{Simple effects} (no verdict).
  \begin{itemize}
  \item UAQ with BF16 rollouts, R2 $-$ R1: late $-1.0$ points (90\% CI $[-3.4, 1.4]$), reward AUC $-0.2$ ($[-1.5, 1.0]$, inside $\pm\delta$), mismatch ratio 0.94.
  \item Global learning rate $\times1.5$, R6 $-$ R3: late $-1.7$ ($[-4.2, 0.9]$), mismatch ratio 0.53 ($[0.23, 1.23]$).
  \end{itemize}
  \end{itemize}
\item \textbf{Update tracking, first look.} All 168 linear weights were re-quantized at each of the 436 syncs in all sixteen INT8 runs, on the CUTLASS kernel. In the INT8 pilot, one PPO step at $10^{-6}$ flipped 0.21\% of layer 0's Q/K/V codes, against 0.20\% predicted by $|d|/\Delta$ (Proposition~4). Row by row the agreement is loose (correlation 0.2--0.4); E2 is the full test.
\end{itemize}

\subsection{E2, static part: U holds at the rounding boundary}\label{sec:e2-results}
\textbf{Data (2026-11-15).} The decoder linear weights of the three stored checkpoints, with the rules of \cref{sec:e2}; the instrumented run and its analyses are in \cref{sec:e2-dyn}. A cell is one projection type of one checkpoint, 21 in all. The 0.5B numbers were seen before registration (\cref{sec:e2}), so only the 1.5B and 7B cells are fresh tests.
\begin{center}\small
\begin{tabularx}{\linewidth}{@{}L l l l l@{}}
\toprule
Rule & 0.5B & 1.5B & 7B & Verdict\\
\midrule
U1: share of $|\varphi|\ge0.45$ over 0.10, in $[0.95,1.05]$ & 1.003--1.011 & 0.999--1.010 & 1.002--1.004 & \textbf{holds}, 21 of 21\\
U2: largest of 40 bin deviations, $|c|\le15$, at most 0.05 & 0.026--0.045 & 0.025--0.052 & 0.015--0.029 & \textbf{fails}, 20 of 21\\
Lemma~1: share of rows with $\sigma/\Delta<1$, below 1\% & 0 & 0 & 0 & \textbf{holds}\\
\bottomrule
\end{tabularx}
\end{center}
\begin{itemize}
\item \textbf{U holds at the rounding boundary.} The near-boundary ratio is 0.999--1.011 in every cell, with the exact ties counted as the rule says. Without them it is 0.985--0.993: the ties sit on the boundary itself, and the band next to it holds slightly less than uniform. Exact ties are 0.18\%, 0.15\% and 0.12\% of the elements at 0.5B, 1.5B and 7B.
\item \textbf{U2 fails narrowly, in one cell:} Qwen2.5-Math-1.5B's \texttt{up\_proj}, at 0.0516. The other 20 cells lie between 0.015 and 0.046, and \texttt{gate\_proj} and \texttt{up\_proj} are the highest at 0.5B and 1.5B. Larger codes deviate more: 0.07--0.13 for $16\le|c|\le63$ and 0.31--0.42 for $64\le|c|\le126$, the BF16 lattice anticipated in \cref{sec:e2}. By the registered reading, a failure of U limits where the propositions apply; it does not refute them.
\item \textbf{Lemma~1's condition holds in every row, by a wide margin.} Median $\sigma/\Delta$ is 28--36, and 99\% of rows have $\sigma/\Delta\ge3.9$; the lowest 1\% quantile is in 7B's \texttt{gate\_proj}, whose outlier rows reach $\max/\rms=32.5$. The theory note's estimate of 6--25 assumed $\max/\rms$ of 5--20; measured, its median is 3.6--4.5 in every cell.
\item \textbf{The link to 7B} (\cref{sec:link7b}): $\nu=\eta/\Delta$ at $\eta=10^{-6}$, reported without a rule.
\begin{center}\small
\begin{tabular}{@{}llll@{}}
\toprule
& 0.5B & 1.5B & 7B\\
\midrule
Median over rows & $1.78\times10^{-3}$ & $8.50\times10^{-4}$ & $1.26\times10^{-3}$\\
10--90\% over rows & $1.21$--$2.27\times10^{-3}$ & $0.66$--$1.02\times10^{-3}$ & $0.96$--$1.51\times10^{-3}$\\
Median by decoder layer & $1.67$--$1.87\times10^{-3}$ & $0.82$--$0.94\times10^{-3}$ & $1.22$--$1.49\times10^{-3}$\\
\bottomrule
\end{tabular}
\end{center}
At the same learning rate, S3's 0.5B runs sit at $1.4\times$ the 7B median, and a 1.5B run would sit at $0.67\times$. Matching the 7B median at 1.5B needs $\eta\approx1.5\times10^{-6}$. By Proposition~4, an update of size $\eta$ flips a code with probability $\nu$: about one code in 560 at 0.5B, 1{,}200 at 1.5B and 790 at 7B.
\end{itemize}

\textbf{Exploratory, not registered: the deviations are those of BF16 storage.} A CPU simulation (\texttt{experiments/e2/bf16\_lattice\_sim.py}) passes synthetic Gaussian rows through the static script's own statistics, once in FP32 and once rounded to BF16 as a checkpoint stores them, with 24M elements per run.
\begin{center}\footnotesize
\begin{tabularx}{\linewidth}{@{}L l l l l l@{}}
\toprule
& & & \multicolumn{3}{c@{}}{Largest bin deviation}\\
\cmidrule(l){4-6}
Rows & U1 & Exact ties & $|c|\le15$ & 16--63 & 64--126\\
\midrule
Synthetic, FP32 & 0.9995 & 0.0002\% & 0.005 & 0.004 & 0.011\\
Synthetic, BF16; row scales spread over 0.75--3 octaves & 1.006--1.009 & 0.15--0.20\% & 0.023--0.032 & 0.073--0.089 & 0.31--0.34\\
Synthetic, BF16; all rows on one scale & 1.014--1.015 & 0.21--0.23\% & 0.070--0.074 & 0.122 & 0.37--0.38\\
Real checkpoints (21 cells) & 0.999--1.011 & 0.07--0.26\% & 0.015--0.052 & 0.070--0.127 & 0.31--0.42\\
\bottomrule
\end{tabularx}
\end{center}
\begin{itemize}
\item \textbf{BF16 rounding alone reproduces the real deviations} in every column. Without it the phases are uniform up to counting noise: 0.005 is the largest of 40 bins, each with a standard deviation of 0.002.
\item \textbf{Other shapes agree.} Student-$t$ rows (5 degrees of freedom) and rows of length 896 or 3584 give the same picture ($|c|\le15$: 0.024--0.027). The simulation does not predict which cell is highest; the failing cell lies between the spread and the equal-scale cases.
\item \textbf{Reading.} BF16 stores the weights on a lattice whose spacing, in grid units, grows with the code, from about $|c|/256$ to $|c|/128$. Lemma~1 assumes a density, so it cannot see this. On the stored checkpoints, U therefore holds at the boundary to within 1.5\%, with or without the ties. In the bulk of small codes the largest bin deviation is 1.5--5.2\%, depending on the cell, which is the size BF16 storage alone produces.
\item \textbf{What it means for the dynamic tests.} Training moves the FP32 master weights off the lattice, so U2 may improve over the run; the dynamic analysis reports U1 and U2 at every full-dump sync. Whether the lattice matters at the scale of a single update, with $|u|$ of order $10^{-3}$ grid steps, far below U1's band of 0.05, is what P4a measures. Its rule leaves out the exact ties in the early pairs. (Written before the run's data; \cref{sec:e2-dyn} has the outcome.)
\end{itemize}

\subsection{E2, dynamic part: the rounding error persists, and the rollout follows the trainer}\label{sec:e2-dyn}
\textbf{The run (2026-11-15).} One run of arm R3 with S3's own call (seed~1, 435 steps) and the dump extension took 4.5\,h on the H100. It behaved like S3's R3 seed~1, the same call without the dump: late accuracy 50.6\% against 50.2\%, and mean $k_3$ $5.9\times10^{-2}$ against $5.5\times10^{-2}$. It had two single-step $k_3$ spikes (0.55 at step 42, 0.60 at step 92), which S3's run did not. The dump check passed at all 14 full-dump syncs: layer~0's live INT8 codes equal the quantizer applied to the dumps, so the dumps are what the rollout quantized. Each full dump has 357.5M elements in the 168 linear weights. The frozen scripts ran after the run ended; none of their output had been seen before.

\textbf{Weight space.} The one-step pairs are 1--2, 2--3, 65--66, 257--258 and 435--436, listed in that order; 87 of 88 checks pass.
\begin{center}\small
\begin{tabularx}{\linewidth}{@{}L l l@{}}
\toprule
Rule (\cref{sec:e2}) & Measured & Verdict\\
\midrule
U1 and U2 at each of the 14 full-dump syncs & 1.002--1.008; U2 0.036 falling to 0.001 & \textbf{hold}, 28 of 28\\
P3, bounded lag, at 13 syncs & at most 0.999997 & \textbf{holds} (exact)\\
P4a, flip share over prediction: without ties; with all elements from sync 65 & 0.912--0.994; 1.002--1.011 & \textbf{holds}, 8 of 8\\
P4b, calibration slope & 0.930, 0.969, 0.987, 0.990, 0.995 & \textbf{holds}, 5 of 5\\
P4c, capture slope, in $[0.95,1.05]$ & 0.931, 0.969, 0.987, 0.990, 0.994 & \textbf{fails} at pair 1--2\\
P6a and P6b at 13 syncs; P6c along 436 syncs & 1.0000--1.0007; at most 0.020; 0.0031 & \textbf{hold}, 27 of 27\\
P7a, error feedback; P7b, stochastic rounding & at most 1.000000; 1.0004 & \textbf{hold}\\
\bottomrule
\end{tabularx}
\end{center}
\begin{itemize}
\item \textbf{U holds throughout, and the bulk grows more uniform.} Sync~1 is the BF16 checkpoint: 0.18\% exact ties and U2 at 0.036, as in the static part. One step later the ties are gone. U2 then falls steadily as training moves the master weights off the BF16 lattice, as \cref{sec:e2-results} expected: 0.031, 0.022, 0.0095, 0.0026 at sync~65 and 0.0010 at sync~436.
\item \textbf{P4 fails at the first pair only.}
  \begin{itemize}
  \item The first update from the stored checkpoint is captured with slope 0.931, against $[0.95,1.05]$. Its flip share (0.912) and calibration slope (0.930) stay inside their wider bands.
  \item U1 held at sync~1 (1.007), so by the registered reading this failure counts against Proposition~4 as stated.
  \item The shortfall shrinks as the weights leave the lattice: capture 0.969 at the second pair and 0.987--0.994 from step 64 on. The exploratory simulation below traces it to BF16 storage.
  \item From step 64 on, the flip share with all elements is within 1.1\% of the prediction.
  \end{itemize}
\item \textbf{Update size} (no rule).
  \begin{itemize}
  \item One PPO step (four Adam updates) moves the codes by $3.6\times10^{-3}$ grid steps on average at the first step, $2.2\times10^{-3}$ at the second and $1.1$--$1.2\times10^{-3}$ from step 64 on. The median $\nu$ is $1.78\times10^{-3}$ (\cref{sec:e2-results}).
  \item Late in training a step flips about 0.12\% of the INT8 codes.
  \item The moving grid adds 2--3\% to the prediction, and phase and update are uncorrelated ($|r|\le0.0006$).
  \end{itemize}
\item \textbf{Proposition~6 holds: the rounding error keeps its size and persists.}
  \begin{itemize}
  \item $\Var(\varphi)$ stays within 0.07\% of its start.
  \item $\Corr(\varphi_k,\varphi_0)$ follows $1-6a(1-a)$ within 0.02 in every bin at every full dump. At sync~436 it runs from 0.87 down to $-0.50$ across the bins.
  \item After 435 steps the rounding error still correlates 0.691 with its start (prediction 0.692).
  \item The grid is nearly fixed, as the proposition assumes. At sync~436 no linear weight is more than $4.8\times10^{-4}$ from its start, so no row's absmax moved by more either. That is 0.7\% of the median row's absmax, and at most 1.3\% for 99\% of rows.
  \end{itemize}
\item \textbf{Proposition~7 holds on the recorded path.} Each scheme is replayed on the path that the run's round-to-nearest rollouts produced; rollouts quantized with error feedback would change the path itself. At $K=436$ syncs the rms partial sum of the errors is:
  \begin{itemize}
  \item $115\,\Delta_0$ for RTN, which is $0.91\,K/\sqrt{12}$, nearly the fully persistent value;
  \item $8.53\,\Delta_0$ for fresh stochastic rounding, against $\sqrt{K/6}=8.52$;
  \item $0.289\,\Delta_0$ for error feedback, which stays at about $1/\sqrt{12}$ for every $K$.
  \end{itemize}
\end{itemize}

\textbf{Exploratory, not registered: the first update starts from BF16 storage.} A CPU simulation (\texttt{experiments/e2/bf16\_update\_sim.py}) passes synthetic rows through the dynamic script's own statistics, stored in BF16 or in FP32, with the same random draws otherwise. Each row's $\nu$ matches the 0.5B checkpoint, and the updates are Adam steps on pure-noise gradients, four per sync. The resulting predicted flip shares are 6--15\% above the run's.
\begin{center}\footnotesize
\begin{tabularx}{\linewidth}{@{}L l l l l l l@{}}
\toprule
& \multicolumn{2}{c}{Pair 1--2} & \multicolumn{2}{c}{Pair 2--3} & \multicolumn{2}{c@{}}{Pair 65--66}\\
\cmidrule(lr){2-3}\cmidrule(lr){4-5}\cmidrule(l){6-7}
Without the sync-1 ties & $F/P$ & Capture & $F/P$ & Capture & $F/P$ & Capture\\
\midrule
E2 run & 0.912 & 0.931 & 0.962 & 0.969 & 0.987 & 0.987\\
Synthetic, BF16 (2 seeds) & 0.913--0.914 & 0.936--0.938 & 0.961--0.963 & 0.968--0.973 & 0.995 & 0.989--0.996\\
Synthetic, FP32 & 0.999 & 1.000 & 1.005 & 1.008 & 0.992 & 0.993\\
\bottomrule
\end{tabularx}
\end{center}
\begin{itemize}
\item \textbf{BF16 storage alone reproduces the shortfall} of the first two pairs, to within 0.008. With FP32 storage, flips match the prediction within twice the counting noise (0.004--0.006). The simulation is too small to resolve the run's 1\% shortfall at pair 65--66 (its noise there is 0.006).
\item \textbf{Reading.} One update moves $q$ by about 0.004 grid steps, 14 times less than the width of U1's band.
  \begin{itemize}
  \item At that scale BF16 storage is visible. Next to a rounding boundary the stored weights are not spread continuously: some sit on it (the ties, which the rule leaves out), and the rest at discrete distances, often larger than one update.
  \item So fewer weights lie within one update of a boundary than U assumes. U fails at the scale of the first update even though U1 holds; the registered U1 test was too coarse to see it.
  \item Training removes the lattice within a few steps.
  \end{itemize}
  This reading explains the registered failure; it does not undo it.
\end{itemize}

\textbf{Policy space.} The probe set is 128 GSM8K training prompts with one response each from $\theta_0$, 30{,}155 tokens in all.
\begin{center}\small
\begin{tabularx}{\linewidth}{@{}L l l@{}}
\toprule
Rule (\cref{sec:e2}) & Measured & Verdict\\
\midrule
P5: slope in $[0.90,1.10]$, cumulative pairs $(1,s')$, $s'\ge9$ & 0.8985--1.0115 (median 0.931) & \textbf{fails}: 1--258 at 0.8985; 9 of 10 hold\\
P5: slope in $[0.90,1.10]$, pooled one-step pairs & 1.169 & \textbf{fails}\\
P5$'$: slope $>0.5$, the same 11 slopes & 0.8985--1.169 & \textbf{holds}, 11 of 11\\
\bottomrule
\end{tabularx}
\end{center}
\begin{itemize}
\item \textbf{The rollout is not frozen.} P5$'$ holds everywhere, so QuRL's claim~(i) finds no support. From sync~1 to 436 the INT8 model's log-probabilities follow the trainer's with slope 0.934 ($r=0.983$); a frozen rollout would give 0.
\item \textbf{Cumulatively it follows at about 93\%.}
  \begin{itemize}
  \item From sync~33 on, $r$ is 0.95--0.98. Six of the eight slopes lie between 0.915 and 0.934; 1--258 is at 0.8985 and 1--435 at 1.011.
  \item So first-order tracking holds within the registered $\pm10\%$ in 9 of 10 pairs, with a consistent shortfall of about 7\% that the theory does not predict.
  \item The slopes' standard errors (0.001--0.004) treat the tokens as independent and are too small. The rules use the slope alone.
  \end{itemize}
\item \textbf{One-step changes are dominated by the rollout's own error, more so late in training.}
  \begin{itemize}
  \item The spread of $\Delta\log\mu$ is 1.7 times that of $\Delta\log\pi$ at the first pair and 8.2 times at the last.
  \item $r$ falls from 0.58 at the first pair to 0.13 and 0.33 at the last two.
  \item The one-step slopes are 1.00, 0.91, 1.03, 0.64 and 2.71. The pooled failure comes from the last pair; without it the pooled slope would be 0.927, a post hoc calculation.
  \end{itemize}
\item \textbf{The mismatch grows while the weight error does not.}
  \begin{itemize}
  \item Mean $|\log\mu-\log\pi|$ over the probe tokens rises from 0.057 at sync~1 to 0.20 at sync~435, 3.5-fold. The weight error keeps its size in grid units (P6a).
  \item The probe contexts are fixed and the weight grids nearly so. The rise must therefore come from growing sensitivity to the weight error or from activation error, not from the size of the weight error (theory note, Section~3.4). The probe responses also become less likely as training proceeds: mean $\log\pi$ falls from $-0.48$ to $-1.35$.
  \item Attributing the rise is E3, left to P4.
  \end{itemize}
\end{itemize}

% =====================================================================================
\section{What the experiments say about the theory}\label{sec:assessment}
\emph{As of 2026-11-15: T, S1, S2, all 22 runs of S3, and E2.} Nothing in the experiments contradicts a theoretical result where its assumptions hold, and the exact results hold wherever they can be tested exactly. One registered E2 check counts against Proposition~4 as stated: the first update from the BF16 checkpoint is captured at 93\%. An exploratory analysis traces it to Assumption~U failing at the scale of one update, below the resolution of U's registered test (\cref{sec:e2-dyn}). S3 adds real-stack evidence that UAQ does not lower the mismatch during training, now also as a registered result (H2's mismatch endpoint), and that the INT8 rollout follows the trainer. Its accuracy tests of Corollary~3 are underpowered: H1, H2 and H3 all came out inconclusive, as predicted. E2 confirms Propositions~6 and~7 on a real run, and Proposition~4 from the second update on. The largest S3 result lies outside the theory's claims: INT8 rollouts cost about 5 points of accuracy here.
\begin{itemize}
\item \textbf{Theorem~1 holds on the real kernel.} Log-probabilities are bit-identical at $s=2$ on vLLM's INT8 kernel at all three sizes (\cref{sec:s1-results}).
  \begin{itemize}
  \item \textbf{No reduction at step 0.} At $s=1.5$ the mismatch ratios are 0.98--1.02, with CIs that exclude an $s$-fold reduction (0.667 and 0.444).
  \item \textbf{No reduction during training.} In S3, R4's mean $k_3$ is $1.04\times$ R3's (90\% CI $[0.87, 1.24]$). Against the BF16 control, H2's registered mismatch test is supported: UAQ's factor on the INT8 mismatch, divided by its factor on the BF16 floor, is 1.10 (90\% CI $[1.015, 1.20]$). The paper's mechanism, error divided by $s$, would give about $1/s^2\approx0.44$.
  \item \textbf{Scope.} R4 against R3 is not a test of Theorem~1 itself, since the trajectories differ through Theorem~2. It rules out, within the margin, a reduction that would appear only later in training.
  \end{itemize}
\item \textbf{Theorem~2 and Corollary~3 hold exactly; RL adds little.} Training is bit-identical in S2's 20-step QuRL loop at both learning rates and in VeRL's FSDP2 actor on the H100 (\cref{sec:s2-results}). Computing the clip norm in the original coordinates, or matching only the learning rates, breaks the equality at step 1, so the theorem's conditions are needed in practice too.
  \begin{itemize}
  \item \textbf{H1 (Remark~3), registered: inconclusive.} Over five pairs R4~$-$~R5 is $+8.5$ points of late accuracy (90\% CI $[-7.5, 24.4]$). Both large values come from R5 stalls; if R4 and R5 stall equally often, both stalls land in R5 with probability 0.22. In the three seeds where neither stalled, R4~$-$~R5 is $-0.9$, $-4.9$ and $+4.1$ points, with no consistent sign.
  \item \textbf{Why H1 stays open.} Nothing contradicts H1, but its test cannot confirm it. Equivalence needed a paired SD below about 2.1 points with five seeds (\cref{sec:prereg}); the stalls made it 16.7, and R4's own seed-to-seed SD is 2.7. This limits the test, not the theorem, which holds bit for bit in S2.
  \item \textbf{UAQ's own effect.} By Corollary~3, R4~$-$~R3 is the effect of the $\times1.5$ Q/K/V and $\div1.5$ norm-gain learning rates. It is small here: late accuracy $-1.1$ points (90\% CI $[-3.3, 1.1]$) and an unchanged mismatch. So the benefit the paper reports for UAQ does not appear in this setting. With BF16 rollouts UAQ's effect is similarly small (late accuracy $-1.0$ points, 90\% CI $[-3.4, 1.4]$).
  \item \textbf{H2 (registered): inconclusive on accuracy, supported on mismatch.} The late-accuracy interaction is $-0.1$ points with a 90\% CI of $[-4.5, 4.3]$, too wide to show equivalence within $\pm\delta$, as the power analysis predicted. One secondary endpoint, final accuracy, detected a small interaction ($+1.7$ points, inside $\pm\delta$): UAQ did slightly better on final accuracy with INT8 rollouts than with BF16. Nothing in the theory predicts such an interaction, and late accuracy shows none, so it is a candidate for a confirmatory run, not a finding.
  \item \textbf{H3 (registered): inconclusive.} By Theorem~2, R6 and R4 are different optimizers: R6 multiplies the learning rate of all 290 actor tensors by 1.5, while UAQ in effect multiplies that of the 72 Q/K/V weights by 1.5 and divides that of the 24 norm gains by 1.5. In accuracy the two cannot be told apart with three seeds (late: $-0.5$ points, 90\% CI $[-5.2, 4.2]$). R6's mismatch was lower than R4's in every seed (ratio 0.51, 90\% CI $[0.26, 1.00]$), which no result here predicts.
  \end{itemize}
\item \textbf{Propositions~3--5: the INT8 rollout follows the trainer.} QuRL's claim~(i) (Section~1.3 of the theory note) predicts a rollout model that stays near its starting point.
  \begin{itemize}
  \item \textbf{Tracking.} In the fourteen INT8 runs of S3 that had learned by step 200, the trainer moved 0.43--0.80 nats per token from the reference model, its starting point, over steps 200--435 (medians). The INT8 rollout stayed 0.031--0.061 from the trainer ($k_3$), 9.7--17.2$\times$ closer. A frozen rollout would show a mismatch of the order of the drift.
  \item \textbf{Learning.} The INT8 model learns: its validation accuracy, generated by the INT8 engine, reaches about 50\%, against 54--57\% with BF16 rollouts. Tracking does not prevent the accuracy gap below.
  \item \textbf{Quantitative test (E2, \cref{sec:e2-dyn}).} On the synced master weights of an instrumented R3 run, Proposition~4 holds from the second update on.
    \begin{itemize}
    \item The capture slope is 0.97--0.99, and from step 64 on the flip share is within 1.1\% of the prediction.
    \item The first update from the BF16 checkpoint is captured at 93\%, a registered failure that BF16 storage reproduces in simulation.
    \item In policy space the rollout follows the trainer with slope 0.93 over the run and never comes near freezing. The registered $\pm10\%$ test fails in 2 of 11 checks: one cumulative pair at 0.8985, and the pooled one-step slope, 1.17, which comes from the last pair.
    \end{itemize}
  \item \textbf{Assumption~U} (\cref{sec:e2-results,sec:e2-dyn}).
    \begin{itemize}
    \item On the stored checkpoints the rounding phase is uniform at the boundary to within 1.5\% in all three models, and Lemma~1's condition holds in every row. In the bulk, BF16 storage leaves deviations of up to 5.2\%, one of them above the registered limit of 5\%.
    \item During the run the bulk deviation falls from 0.036 to 0.001 as the weights leave the BF16 lattice.
    \item At the scale of one update the lattice matters, for the first update only.
    \end{itemize}
  \end{itemize}
\item \textbf{Propositions~6--7 hold on a real run.} In E2 (\cref{sec:e2-dyn}), the rounding error persists as Proposition~6 predicts, within 0.02 in every bin.
  \begin{itemize}
  \item After 435 steps it still correlates 0.69 with its start, and it keeps its size.
  \item Replayed on the recorded path, error feedback bounds the partial sums, while RTN's grow almost linearly (Proposition~7).
  \end{itemize}
  S3 sharpens the question Proposition~6 raises.
  \begin{itemize}
  \item \textbf{The mismatch rise is unexplained.} In the ten H1 runs without a stall, $k_3$ rises 4.7--7.8-fold, from $1.0$--$1.2\times10^{-2}$ at step 1 to window medians of $5.5$--$8.4\times10^{-2}$ over steps 60--175, then partly falls.
    \begin{itemize}
    \item E2 confirms that the weight error keeps its size while the mismatch on a fixed probe set rises 3.5-fold. In its run the weight grids moved by at most 1.3\% (99\% of rows).
    \item So the rise must come from sensitivity, activation error or contexts, not from growing weight grids, at least in E2's run.
    \item It tracks neither entropy nor the drift from the reference model. Attributing it is E3.
    \end{itemize}
  \item \textbf{The side hypothesis of the theory note's Section~3.5 is not supported.} UAQ's slower norm gains were to slow mismatch growth. R4's mismatch is not lower than R3's, either on average ($1.04\times$) or at the end of training (steps 400--435: $1.13\times$, $1.56\times$ and $0.88\times$ per seed).
  \item \textbf{What is at stake.} INT8 rollouts cost about 5 points of accuracy here: 4.9 points averaged over UAQ (95\% CI $[2.3, 7.6]$), and every BF16 run ends above every INT8 run. The gap comes with lower training reward, higher entropy and more drift from the reference model. Several causes are open:
  \begin{itemize}
  \item the persistent rounding error of Proposition~6;
  \item the bias of TIS's truncation;
  \item scoring on the INT8 engine.
  \end{itemize}
  Separating them is E3's task. Closing the gap is what P6 (dithered rollouts) would have to show.
  \end{itemize}
\item \textbf{Outside the theory.} The early stalls (4 of 22 runs) and the failure at $10^{-5}$ are optimization behaviour of this PPO configuration.
  \begin{itemize}
  \item \textbf{Stalls and quantization.} All four stalls had INT8 rollouts, against none of the six BF16 runs, but with $p=0.25$ this does not tie them to quantization. In seed~3, with the same prompt order, three of the four INT8 arms stalled and neither BF16 run did.
  \item \textbf{Not the mismatch level.} Within the INT8 runs, the stalls' mismatch ranged from the highest of any run (R3 seed~3) to among the lowest (R5's two, R6 seed~3), so the mismatch level does not explain them.
  \item \textbf{Dips.} The weak link between validation dips and mismatch appears in the BF16 runs too ($+0.18$), so it is not about quantization. INT8 validation curves are bumpier, though.
  \end{itemize}
\item \textbf{For the write-up.} The case that UAQ is an optimizer reparameterization (chiefly per-layer learning rates), not a quantization fix, rests on exact identities verified on the real stack (\cref{sec:s1-results,sec:s2-results}) and on the absence of a mismatch reduction, both at step 0 and during RL. The RL accuracy results are consistent with it but underpowered: H1, H2 and H3 came out inconclusive on accuracy, as predicted, while H2's mismatch endpoint, registered in advance, supports it. E2 shows on a real run that the INT8 rollout tracks the trainer and that its rounding error persists, as the theory predicts. The largest remaining gap is E3 (what drives the mismatch rise), left to P4. Separately, QuRL's premise that INT8 rollouts cost little accuracy holds less well here than in the paper: about 5 points against 1.8 (95\% CI $[2.3, 7.6]$ averaged over UAQ, three seeds). That is the effect E3 and P6 would have to explain and reduce.
\end{itemize}

% =====================================================================================
\section{Status}\label{sec:status}
{\small
\begin{xltabular}{\linewidth}{@{}L l L L@{}}
\toprule
Claim & Theory & Evidence so far & Status\\
\midrule
\endhead
UAQ leaves the rollout model unchanged & Thm.~1 & float64; vLLM INT8 bit-identical at $s=2$ (3 sizes); no change at $s=1.5$; in RL, R4's mean $k_3$ within 0.80--1.25 of R3's, and H2's registered mismatch test supported (ratio of ratios 1.10) & established\\
UAQ $=$ matched optimizer settings & Thm.~2, Cor.~3 & float64; bit-identical 20-step loop; VeRL check; H1 (registered, 5 seeds): R4 $-$ R5 $=+8.5$ points late, 90\% CI $[-7.5, 24.4]$ & exact equivalence established; H1 in RL inconclusive\\
Rounding phase is uniform & Assum.~U, Lem.~1 & E2 static, 21 cells (0.5B, 1.5B, 7B): boundary share within 1.1\% of uniform; bulk ($|c|\le15$) within 5\% in 20 (one at 5.2\%, BF16 storage); $\sigma/\Delta\ge1$ in every row. E2 run: U1, U2 hold at all 14 syncs & holds as registered but for one bulk cell; BF16 storage breaks it at the scale of one update\\
Updates are captured & Props.~3--5 & float64; S3 pilot 0.21\% vs.\ 0.20\%. E2: capture slope 0.97--0.99 after the first update, 0.93 at it; policy-space slope 0.93, not frozen & established after the first update; P5's $\pm10\%$ band failed in 2 of 11\\
Error is persistent; error feedback bounds it & Props.~6--7 & float64; E2: $\Corr(\varphi_k,\varphi_0)$ within 0.02 of $1-6a(1-a)$, 0.69 after 435 steps; error feedback bounded & established on a real run; P6 pending\\
UAQ's effect is not quantization-specific & Cor.~3 & H2 (registered, 3 seeds): late accuracy $-0.1$ points, 90\% CI $[-4.5, 4.3]$; mismatch ratio of ratios 1.10, $[1.015, 1.20]$; final accuracy $+1.7$ (secondary) & accuracy inconclusive; mismatch supported\\
A global learning rate is the wrong control & Thm.~2 & H3 (registered, 3 seeds): R6 $-$ R4 $=-0.5$ points late, $[-5.2, 4.2]$; inconclusive on every endpoint & inconclusive\\
QuRL's Table~1 at $\eta=10^{-5}$ & --- & not reproduced & open (ask the authors)\\
INT8 rollouts cost little accuracy (Table~1: 1.8 points) & --- & at $\eta=10^{-6}$: BF16 53.9--56.6\% late (6 runs), INT8 at most 53.0\% (16 runs); gap 4.9 points averaged over UAQ (95\% CI $[2.3, 7.6]$); INT8 scored on the INT8 engine & gap larger than the paper's; cause open (E3, P6)\\
\bottomrule
\end{xltabular}}


% Background references listed in the bibliography but not cited in the text above.
\nocite{yao2025flashrl,hilton2022batch}

\bibliographystyle{plainnat}
\bibliography{refs}

\end{document}
